% Source: physical PDF pp. 98--106 (printed pp. 75--83), from the Section 26.4
% heading through the final paragraph immediately before Section 26.5.
% Inventory: Eqs. (26.4.1)--(26.4.21), nine unnumbered display groups, one
% footnote, citation [2], and no figures, tables, or typographic dividers.
% Source oddity preserved: Eq. (26.4.15) contains a repeated inner sum over n.
% Uncertainties: none.

\subsection{Renormalizable Theories of Chiral Superfields}

\providecommand{\SuperQ}{\mathcal{Q}}

We will now work out the details of a general renormalizable theory of scalar
chiral superfields. This will provide some insight into the implications of
supersymmetry, and the theory we obtain will be part of the supersymmetric
standard model to be discussed in Chapter 28.

As discussed in Section 12.2, the Lagrangian density in a renormalizable theory
can contain only operators with dimensionality (counting powers of energy or
momentum, with \(\hbar=c=1\)) four or less. Eq.~\eqref{eq:26.2.6} shows that
\(\SuperQ_\alpha\) and hence \(\partial/\partial\theta^\alpha\) has
dimensionality \(1/2\), so \(D_\alpha\) and \(\bar D_{\dot\alpha}\) have
dimensionality \(+1/2\), and
\(\theta_\alpha\) has dimensionality \(-1/2\). The \(\mathcal{F}\)- and
\(D\)-terms of a superfield \(S\) are the coefficients of two or four factors
of \(\theta\), respectively, so if the superfield has dimensionality \(d(S)\)
then its \(\mathcal{F}\)- and \(D\)-terms have dimensionalities
\(d(\mathcal{F}^S)=d(S)+1\) and \(d(D^S)=d(S)+2\). Thus in a renormalizable
theory the functions \(f\) and \(K\) in Eq.~\eqref{eq:26.3.30} consist of
operators with dimensionality at most three and two, respectively.

The dimensionality of an elementary scalar superfield \(\Phi_n\) is that of an
elementary scalar field, or \(+1\), so in order for each term in the function
\(f\) to have dimensionality three or less, it can contain at most three
factors of \(\Phi_n\) and/or derivatives \(\partial/\partial x^\mu\) and/or
pairs of spinor superderivatives \(D_\alpha,\bar D_{\dot\alpha}\). As discussed in the
previous section, any left-chiral term in \(f\) that involves superderivatives
could be replaced with a term in \(K\), so superderivatives can be omitted in
\(f\). Eq.~\eqref{eq:26.2.30} shows that spacetime derivatives can be expressed
in terms of superderivatives, so these too can be omitted. (In any case,
Lorentz invariance would rule out terms with one spacetime derivative, and
terms with two derivatives in a renormalizable theory could involve only one
\(\Phi_n\) factor, on which these derivatives would have to act, so such terms
would not contribute to the action.) We conclude that \(f(\Phi)\) is at most a
cubic polynomial in the \(\Phi_n\), without spacetime derivatives or
superderivatives.

The same dimensional analysis shows that in a renormalizable theory \(K\) is at
most a quadratic function of \(\Phi_n\) and \(\Phi_n^*\), without derivatives.
But any term in \(K(\Phi,\Phi^*)\) that involves only \(\Phi_n\) or only
\(\Phi_n^*\) would be a chiral superfield, and chiral superfields by definition
have no \(D\)-terms, so \([K(\Phi,\Phi^*)]_D\) receives contributions only from
terms in \(K(\Phi,\Phi^*)\) that involve both \(\Phi_n\) and \(\Phi_n^*\). Thus
\(K(\Phi,\Phi^*)\) must be of the form
\begin{equation}
  K(\Phi,\Phi^*)=\sum_{mn}g_{nm}\Phi_n^*\Phi_m\,,
  \tag{26.4.1}\label{eq:26.4.1}
\end{equation}
with constant coefficients \(g_{nm}\) forming a Hermitian matrix.

We must now calculate the \(\mathcal{F}\)- and \(D\)-components of \(f(\Phi)\)
and \(K(\Phi,\Phi^*)\), respectively. To calculate the \(D\)-component of
\(K(\Phi,\Phi^*)\), we note that the term in \(\Phi_n^*\Phi_m\) of fourth order
in \(\theta\) is
\[
  \begin{aligned}
  [\Phi_n^\dagger\Phi_m]_{\theta^2\bar\theta^2}
  =\theta\theta\bar\theta\bar\theta\biggl[
    &\frac14\phi_n^*\Box\phi_m
    +\frac14(\Box\phi_n^*)\phi_m
    -\frac12\partial_\mu\phi_n^*\partial^\mu\phi_m\\
    &+\frac{i}{2}\bar\psi_n\bar\sigma^\mu\partial_\mu\psi_m
    -\frac{i}{2}(\partial_\mu\bar\psi_n)\bar\sigma^\mu\psi_m
    +\mathcal F_n^*\mathcal F_m\biggr].
  \end{aligned}
\]
The two-spinor identities \eqref{eq:26.A.18} and
\eqref{eq:26.A.19} put all quartic monomials into the single standard order:
\[
  (\theta\sigma^\mu\bar\theta)
  (\theta\sigma^\nu\bar\theta)
  =-\frac12\theta\theta\bar\theta\bar\theta\,\eta^{\mu\nu},
  \qquad
  (\theta\psi_n)(\bar\theta\bar\psi_m)
  =-\frac12(\theta\sigma^\mu\bar\theta)
    (\psi_n\sigma_\mu\bar\psi_m).
\]
The \(D\)-term is twice the coefficient of
\(\theta\theta\bar\theta\bar\theta\), minus \(\frac12\Box\) acting on the
\(\theta\)-independent term, which for \(\Phi_n^\dagger\Phi_m\) is
\(\phi_n^*\phi_m\), so
\begin{equation}
  \begin{aligned}
  \frac12[K(\Phi,\Phi^*)]_D
  =\sum_{nm}g_{nm}\biggl[
    &-\partial_\mu\phi_n^*\partial^\mu\phi_m
    +\mathcal{F}_n^*\mathcal{F}_m
  \\
    &+i\bar\psi_n\bar\sigma^\mu\partial_\mu\psi_m
  \biggr]\,.
  \end{aligned}
  \tag{26.4.2}\label{eq:26.4.2}
\end{equation}
If we write \(\Phi_n\) as a linear combination
\(\sum_mN_{nm}\Phi'_m\) of new superfields \(\Phi'_m\), then
\(K(\Phi,\Phi^*)\) is given in terms of the new superfields by a formula which
is the same as Eq.~\eqref{eq:26.4.1}, except that \(g_{nm}\) is replaced with
\(g'_{nm}=(N^\dagger gN)_{nm}\). In order for the kinematic terms for the
scalar and spinor fields to have a sign consistent with the quantum
commutation and anticommutation relations, it is necessary that the Hermitian
matrix \(g_{nm}\) be positive-definite and, as shown in Section 12.5, this
means that we can choose \(N\) so that \(g'_{nm}=\delta_{nm}\). Dropping
primes, the term \eqref{eq:26.4.2} is now
\begin{equation}
  \begin{aligned}
  \frac12[K(\Phi,\Phi^*)]_D
  =\sum_n\biggl[
    &-\partial_\mu\phi_n^*\partial^\mu\phi_n
    +\mathcal{F}_n^*\mathcal{F}_n
  \\
    &+i\bar\psi_n\bar\sigma^\mu\partial_\mu\psi_n
  \biggr]\,.
  \end{aligned}
  \tag{26.4.3}\label{eq:26.4.3}
\end{equation}
We can still use a unitary transformation to redefine the superfields without
changing the form of Eq.~\eqref{eq:26.4.3}, a freedom we will need to exercise
shortly.

The terms in Eq.~\eqref{eq:26.4.3} involving \(\phi_n\) and \(\psi_n\) are
the correct kinematic Lagrangians for conventionally normalized complex scalar
and Weyl fields.

To calculate the \(\mathcal{F}\)-term of \(f(\Phi)\), it is most convenient to
use the representation \eqref{eq:26.3.21} of the superfields, and pick out the
term of second order in \(\theta\):
\[
  \begin{aligned}
  [f(\Phi(x,\theta,\bar\theta))]_{\theta^2}
  ={}&\sum_{nm}
    (\theta\psi_n)(\theta\psi_m)
    \frac{\partial^2f(\phi(x))}
      {\partial\phi_n(x)\partial\phi_m(x)}
  \\
  &+\sum_n\mathcal{F}_n(x)
    \frac{\partial f(\phi(x))}{\partial\phi_n(x)}
    \theta\theta\,.
  \end{aligned}
\]
(We have replaced \(x_+\) here with \(x\) because the term
\(\theta\sigma^\mu\bar\theta\) in Eq.~\eqref{eq:26.3.21}
vanishes when multiplied by an expression with two \(\theta\) factors.)
The \(\theta\) dependence of the first term on the right-hand side can be put
into standard form using Eq.~\eqref{eq:26.A.11}:\footnote{The order of the
Grassmann spinors in this identity is part of the convention; no cosmetic
interchange has been made.}
\[
  (\theta\psi_n)(\theta\psi_m)
  =-\frac12(\psi_n\psi_m)(\theta\theta)\,.
\]
The \(\mathcal{F}\)-term of any left-chiral superfield is the coefficient of
\(\theta\theta\), so here
\begin{equation}
  [f(\Phi)]_{\mathcal{F}}
  =-\frac12\sum_{nm}
    \frac{\partial^2f(\phi_n)}{\partial\phi_n\partial\phi_m}
    (\psi_n\psi_m)
  +\sum_n\mathcal{F}_n\frac{\partial f(\phi)}{\partial\phi_n}\,.
  \tag{26.4.4}\label{eq:26.4.4}
\end{equation}
The complete Lagrangian density is the sum of the terms
\eqref{eq:26.4.3}, \eqref{eq:26.4.4}, and the complex conjugate of
\eqref{eq:26.4.4}:
\begin{equation}
  \begin{aligned}
  \mathcal{L}={}&
  \sum_n\biggl[
    -\partial_\mu\phi_n^*\partial^\mu\phi_n
    +\mathcal{F}_n^*\mathcal{F}_n
    +i\bar\psi_n\bar\sigma^\mu\partial_\mu\psi_n
  \biggr]
  \\
  &-\frac12\sum_{nm}
    \frac{\partial^2f(\phi)}{\partial\phi_n\partial\phi_m}
    (\psi_n\psi_m)
  -\frac12\sum_{nm}
    \left(\frac{\partial^2f(\phi)}
      {\partial\phi_n\partial\phi_m}\right)^*
    (\bar\psi_n\bar\psi_m)
  \\
  &+\sum_n\mathcal{F}_n\frac{\partial f(\phi)}{\partial\phi_n}
  +\sum_n\mathcal{F}_n^*
    \left(\frac{\partial f(\phi)}{\partial\phi_n}\right)^* .
  \end{aligned}
  \tag{26.4.5}\label{eq:26.4.5}
\end{equation}
The auxiliary fields \(\mathcal{F}_n\) enter quadratically in the action, with
constant coefficients for the second-order terms, so they can be eliminated
by setting \(\mathcal{F}_n\) equal to the value at which the Lagrangian density
\eqref{eq:26.4.5} is stationary with respect to \(\mathcal{F}_n\) and
\(\mathcal{F}_n^*\):
\begin{equation}
  \mathcal{F}_n=-\left(\frac{\partial f(\phi)}{\partial\phi_n}\right)^* .
  \tag{26.4.6}\label{eq:26.4.6}
\end{equation}
Inserting this in Eq.~\eqref{eq:26.4.5} gives
\begin{equation}
  \begin{aligned}
  \mathcal{L}={}&
  \sum_n\biggl[
    -\partial_\mu\phi_n^*\partial^\mu\phi_n
    +i\bar\psi_n\bar\sigma^\mu\partial_\mu\psi_n
  \biggr]
  \\
  &-\frac12\sum_{nm}
    \frac{\partial^2f(\phi)}{\partial\phi_n\partial\phi_m}
    (\psi_n\psi_m)
  -\frac12\sum_{nm}
    \left(\frac{\partial^2f(\phi)}
      {\partial\phi_n\partial\phi_m}\right)^*
    (\bar\psi_n\bar\psi_m)
  \\
  &-\sum_n
    \left(\frac{\partial f(\phi)}{\partial\phi_n}\right)^*
    \frac{\partial f(\phi)}{\partial\phi_n}\,.
  \end{aligned}
  \tag{26.4.7}\label{eq:26.4.7}
\end{equation}
Thus the scalar field potential is
\(V(\phi)=\sum_n|\partial f(\phi)/\partial\phi_n|^2\).

With the auxiliary fields eliminated in this way, the action is no longer
invariant under the supersymmetry transformations \eqref{eq:26.3.15},
\eqref{eq:26.3.17} of the remaining fields \(\psi_{n\alpha}\) and \(\phi_n\):
\[
  \delta\psi_{n\alpha}
  =-i\sqrt2\,\sigma^\mu_{\alpha\dot\alpha}
    \bar\xi^{\dot\alpha}\partial_\mu\phi_n
  -\sqrt2\left(\frac{\partial f(\phi)}
    {\partial\phi_n}\right)^*\xi_\alpha\,,
  \qquad
  \delta\phi_n=\sqrt2\xi\psi_n\,.
\]
This is because the expression \eqref{eq:26.4.6} does not obey the
transformation rule
\(\delta\mathcal{F}_n
=-i\sqrt2\bar\xi\bar\sigma^\mu\partial_\mu\psi_n\) for
\(\mathcal{F}_n\) given by Eq.~\eqref{eq:26.3.16}, but instead
\[
  \begin{aligned}
  \delta\left(-\frac{\partial f(\phi)}{\partial\phi_n}\right)^*
  &=-\sum_m
    \left(\frac{\partial^2f(\phi)}
      {\partial\phi_n\partial\phi_m}\right)^*
    \delta\phi_m^*
  \\
  &=-\sqrt2\sum_m
    \left(\frac{\partial^2f(\phi)}
      {\partial\phi_n\partial\phi_m}\right)^*
    (\bar\xi\bar\psi_m)\,.
  \end{aligned}
\]
For the same reason, after eliminating the auxiliary fields the commutators of
the supersymmetry transformations of \(\phi_n\) and \(\psi_{n\alpha}\) are no
longer given by the supersymmetry anticommutation relations, and in fact do
not form a closed Lie superalgebra. But this is not inconsistent with the
existence of quantum mechanical operators \(Q_\alpha\) that satisfy the
anticommutation relations of supersymmetry. These operators generate
supersymmetry transformations, in the sense that the commutator of
\(-i(\xi Q+\bar\xi\bar Q)\) with any Heisenberg-picture quantum field
\(\phi_n\) or \(\psi_{n\alpha}\) equals its change under a supersymmetry
transformation with infinitesimal parameter \(\xi\). With
\(\mathcal{F}_n\) given by Eq.~\eqref{eq:26.4.6} the commutator of
\(-i(\xi Q+\bar\xi\bar Q)\) with \(\mathcal{F}_n\) is given by
\(\delta\mathcal{F}_n
=-i\sqrt2\bar\xi\bar\sigma^\mu\partial_\mu\psi_n\), because in
the Heisenberg picture the Weyl field \(\psi_n\) satisfies the field
equation derived from the Lagrangian \eqref{eq:26.4.7}:
\[
  i\bar\sigma^\mu\partial_\mu\psi_n
  =\sum_m
    \left(\frac{\partial^2f(\phi)}
      {\partial\phi_n\partial\phi_m}\right)^*
    \bar\psi_m\,.
\]
Likewise, the supersymmetry transformations of the quantum fields \(\phi_n\)
and \(\psi_{n\alpha}\) do form a closed Lie superalgebra when the field equations
are taken into account. Such algebras are often called on-shell.

The zeroth-order expectation values \(\phi_{n0}\) of the scalar fields
\(\phi_n\) must be at the maximum of the last term in Eq.~\eqref{eq:26.4.7}.
Since this term is always negative or zero, the maximum will be at
spacetime-independent field values \(\phi_{n0}\) at which this term vanishes,
so that
\begin{equation}
  \left.
  \frac{\partial f(\phi)}{\partial\phi_n}
  \right|_{\phi=\phi_0}=0\,,
  \tag{26.4.8}\label{eq:26.4.8}
\end{equation}
provided of course that a solution of this equation exists.
Eq.~\eqref{eq:26.4.8} not only maximizes the last term in
Eq.~\eqref{eq:26.4.7}---it is also the condition for supersymmetry to be
unbroken. Invariance of the vacuum under supersymmetry transformations
requires that the vacuum expectation value of the change in any field under a
supersymmetry transformation should vanish. The change in a bosonic field is
a fermionic field, which of course has zero expectation value anyway, but
Eq.~\eqref{eq:26.3.15} shows that the vacuum expectation value of
\(\delta\psi_{n\alpha}\) is proportional to the vacuum expectation value of the
auxiliary field \(\mathcal{F}_n\), which therefore must vanish if
supersymmetry is unbroken. According to Eq.~\eqref{eq:26.4.6}, in zeroth-order
perturbation theory this condition requires that Eq.~\eqref{eq:26.4.8} must be
satisfied. We will see in Section 27.6 that if Eq.~\eqref{eq:26.4.8} is
satisfied then supersymmetry is unbroken to all orders of perturbation theory.

For a single left-chiral scalar superfield \(\Phi\), the fundamental theorem
of algebra tells us that the polynomial \(\partial f(\phi)/\partial\phi\)
always has at least one zero somewhere in the complex plane. This is not
necessarily true when there is more than one superfield. If we assume that
there is a solution \(\phi_{n0}\) of Eq.~\eqref{eq:26.4.8}, we can evaluate the
physical degrees of freedom of the theory by setting
\begin{equation}
  \phi_n=\phi_{n0}+\varphi_n\,,
  \tag{26.4.9}\label{eq:26.4.9}
\end{equation}
and expanding in powers of \(\varphi_n\). The masses of the particles of this
theory can be calculated by inspecting the terms of second order in
\(\varphi\) and \(\psi\):
\begin{equation}
  \begin{aligned}
  \mathcal{L}_0={}&
  \sum_n\biggl[
    -\partial_\mu\varphi_n^*\partial^\mu\varphi_n
    +i\bar\psi_n\bar\sigma^\mu\partial_\mu\psi_n
  \biggr]
  \\
  &-\frac12\sum_{nm}\mathcal{M}_{nm}(\psi_n\psi_m)
  -\frac12\sum_{nm}\mathcal{M}_{nm}^*
    (\bar\psi_n\bar\psi_m)
  -\sum_{nm}(\mathcal{M}^\dagger\mathcal{M})_{mn}
    \varphi_m^*\varphi_n\,,
  \end{aligned}
  \tag{26.4.10}\label{eq:26.4.10}
\end{equation}
where \(\mathcal{M}\) is the symmetric complex matrix
\begin{equation}
  \mathcal{M}_{mn}\equiv
  \left(
    \frac{\partial^2f(\phi)}{\partial\phi_n\partial\phi_m}
  \right)_{\phi=\phi_0}.
  \tag{26.4.11}\label{eq:26.4.11}
\end{equation}
Now if we redefine the fields by a unitary transformation
\begin{equation}
  \varphi_n=\sum_m\mathcal{U}_{nm}\varphi'_m\,,
  \qquad
  \psi_{n\alpha}=\sum_m\mathcal{U}_{nm}\psi'_{m\alpha}\,,
  \tag{26.4.12}\label{eq:26.4.12}
\end{equation}
then the free-field Lagrangian \eqref{eq:26.4.10} will take the same form, but
with \(\mathcal{M}\) replaced with \(\mathcal{M}'\), where
\begin{equation}
  \mathcal{M}'=\mathcal{U}^{\mathsf T}\mathcal{M}\mathcal{U}\,.
  \tag{26.4.13}\label{eq:26.4.13}
\end{equation}
According to a theorem of matrix algebra, for any complex symmetric matrix
\(\mathcal{M}\) it is always possible to find a unitary matrix \(\mathcal{U}\)
such that the matrix \(\mathcal{M}'\) defined by Eq.~\eqref{eq:26.4.13} is
diagonal, with real positive elements \(m_n\) on the diagonal. (For future
use, we note that
\(\mathcal{M}'^\dagger\mathcal{M}'=
\mathcal{U}^\dagger\mathcal{M}^\dagger\mathcal{M}\mathcal{U}\), so the
quantities \(m_n^2\) are just the eigenvalues of the positive Hermitian matrix
\(\mathcal{M}^\dagger\mathcal{M}\).) Redefining the fields in this way and
dropping the primes, the quadratic part of the Lagrangian is now
\begin{equation}
  \begin{aligned}
  \mathcal{L}_0={}&
  \sum_n\biggl[
    -\partial_\mu\varphi_n^*\partial^\mu\varphi_n
    +i\bar\psi_n\bar\sigma^\mu\partial_\mu\psi_n
  \biggr]
  \\
  &-\frac12\sum_n m_n(\psi_n\psi_n+\bar\psi_n\bar\psi_n)
  -\sum_n m_n^2\varphi_n^*\varphi_n\,.
  \end{aligned}
  \tag{26.4.14}\label{eq:26.4.14}
\end{equation}
The Weyl kinetic term is already in its standard form.  Using the exchange
property \eqref{eq:26.A.7}, the symmetric form
obtained directly from the comparison dictionary differs only by a
derivative:
\[
  \frac{i}{2}\left(
    \psi_n\sigma^\mu\partial_\mu\bar\psi_n
    +\bar\psi_n\bar\sigma^\mu\partial_\mu\psi_n\right)
  =i\bar\psi_n\bar\sigma^\mu\partial_\mu\psi_n
  +\frac{i}{2}\partial_\mu(\psi_n\sigma^\mu\bar\psi_n).
\]
The reality property \eqref{eq:26.A.21} shows that the real mass term is the
sum of the undotted contraction and its conjugate:
\[
  \psi_n\psi_n+\bar\psi_n\bar\psi_n
  =2\operatorname{Re}(\psi_n\psi_n).
\]
The complete quadratic Lagrangian is then
\begin{equation}
  \mathcal{L}_0
  =\sum_n\left[
    -\partial_\mu\varphi_n^*\partial^\mu\varphi_n
    -\sum_n m_n^2\varphi_n^*\varphi_n
    +i\bar\psi_n\bar\sigma^\mu\partial_\mu\psi_n
    -\frac{m_n}{2}(\psi_n\psi_n+\bar\psi_n\bar\psi_n)
  \right].
  \tag{26.4.15}\label{eq:26.4.15}
\end{equation}
The factor \(1/2\) avoids double-counting a Weyl mass term and its conjugate,
while there is no factor \(1/2\) in the scalar terms because these are
complex scalars. We see that the spinless and spin \(1/2\) particles
have equal masses \(m_n\), as required by the unbroken supersymmetry of the
theory.

The interaction part \(\mathcal{L}'\) of the Lagrangian density is given by the
terms in Eq.~\eqref{eq:26.4.7} of higher than second order in \(\varphi_n\) and
\(\psi_n\). Since the superpotential \(f(\phi_0+\varphi)\) is being assumed to
be a cubic polynomial and stationary at \(\varphi_n=0\), with \(\varphi\)
defined so that the second-order terms are
\(\frac12\sum_n m_n\varphi_n^2\), we may write the superpotential (apart from
an inconsequential constant term) as
\begin{equation}
  f(\phi_0+\varphi)
  =\frac12\sum_n m_n\varphi_n^2
  +\frac16\sum_{nm\ell}f_{nm\ell}\varphi_n\varphi_m\varphi_\ell\,.
  \tag{26.4.16}\label{eq:26.4.16}
\end{equation}
Using this in Eq.~\eqref{eq:26.4.7} gives the interaction as
\begin{equation}
  \begin{aligned}
  \mathcal{L}'={}&
  -\frac12\sum_{nm\ell}f_{nm\ell}\varphi_n
    (\psi_m\psi_\ell)
  \\
  &-\frac12\sum_{nm\ell}f_{nm\ell}^*\varphi_n^*
    (\bar\psi_m\bar\psi_\ell)
  \\
  &-\frac12\sum_{nm\ell}m_nf_{nm\ell}
    \varphi_n^*\varphi_m\varphi_\ell
  -\frac12\sum_{nm\ell}m_nf_{nm\ell}^*
    \varphi_n\varphi_m^*\varphi_\ell^*
  \\
  &-\frac14\sum_{nm\ell m'\ell'}
    f_{nm\ell}f_{nm'\ell'}^*
    \varphi_m\varphi_\ell\varphi_{m'}^*\varphi_{\ell'}^*\,.
  \end{aligned}
  \tag{26.4.17}\label{eq:26.4.17}
\end{equation}
We see that a knowledge of the masses \(m_n\) and the `Yukawa' couplings
\(f_{nm\ell}\) of the scalars and fermions is enough to determine all the cubic
and quartic self-couplings of the spinless fields.

As an illustration, consider the case of a single left-chiral superfield. For
comparison with earlier results, let us write the single coefficient \(f\) in
Eq.~\eqref{eq:26.4.16} as
\begin{equation}
  f\equiv2\sqrt2\,e^{i\alpha}\lambda\,,
  \tag{26.4.18}\label{eq:26.4.18}
\end{equation}
where \(\lambda\) is real and \(\alpha\) is some real phase. We will also
introduce a pair of real spinless fields \(A(x)\) and \(B(x)\) by writing the
single complex scalar here as
\begin{equation}
  \varphi\equiv e^{-i\alpha}\left(\frac{A+iB}{\sqrt2}\right).
  \tag{26.4.19}\label{eq:26.4.19}
\end{equation}
The total Lagrangian density is then given by the sum of
Eqs.~\eqref{eq:26.4.15} and \eqref{eq:26.4.17} as
\begin{equation}
  \begin{aligned}
  \mathcal{L}={}&
  -\frac12\partial_\mu A\,\partial^\mu A
  -\frac12\partial_\mu B\,\partial^\mu B
  -\frac12m^2(A^2+B^2)
  \\
  &+\frac{i}{2}\left(
    \psi\sigma^\mu\partial_\mu\bar\psi
    +\bar\psi\bar\sigma^\mu\partial_\mu\psi\right)
  -\frac12m(\psi\psi+\bar\psi\bar\psi)
  -\lambda A(\psi\psi+\bar\psi\bar\psi)
  -i\lambda B(\psi\psi-\bar\psi\bar\psi)
  \\
  &-m\lambda A(A^2+B^2)
  -\frac12\lambda^2(A^2+B^2)^2\,.
  \end{aligned}
  \tag{26.4.20}\label{eq:26.4.20}
\end{equation}
which is the same as the Lagrangian density \eqref{eq:24.2.9} originally found
by Wess and Zumino~\hyperref[ch26-ref-2]{[2]}. It is noteworthy that in this
simple case the Lagrangian turns out to be invariant under a space inversion
transformation
\begin{equation}
  A(x)\longrightarrow A(\Lambda_Px)\,,
  \qquad
  B(x)\longrightarrow-B(\Lambda_Px)\,,
  \qquad
  \psi_\alpha(x)\longrightarrow
  i(\sigma^0)_{\alpha\dot\alpha}
    \bar\psi^{\dot\alpha}(\Lambda_Px)\,,
  \tag{26.4.21}\label{eq:26.4.21}
\end{equation}
even though we did not assume parity conservation in deriving it. The
appearance of parity conservation as an `accidental' symmetry is a familiar
feature of various renormalizable gauge theories (see Sections 12.5 and 18.7)
but not of theories involving spinless fields, so this is a special
consequence of supersymmetry in the renormalizable theory of a single scalar
superfield.
